\section{Scene-graph ingest and merge policy}
\label{app:graph_fusion_policy}

Dynagraph's scene graph is a single shared index used for both the EQA prompt
(SCENE\_GRAPH) and count/FIND recall. This appendix records how detections enter
that graph, why a naive instance feed can flood it, and the merge policy we use to
keep one graph instead of two.

\subsection{Instance-graph ingest and the singleton flood}
\label{app:instance_flood}

Per-frame YoloE instance masks are the richest per-view signal: each instance yields a
median world XYZ, a 3D bounds box, a mask bbox, and a class string
(\texttt{frame\_instances\_to\_detections}). Feeding every instance of every frame
into the scene graph as a first-class node triples-to-quadruples the object-node count
per episode (e.g.\ q8 goes from 19 object nodes to 241 at an unchanged VLM budget of
$\sim$57 planning steps) and drives the singleton fraction (support count of one) to
40--53\%. We reproduced this in the 2026-09-02 full-113 run: 16 of 81 overlap qids
regressed versus the pre-instance-graph baseline, almost all location/state questions
(``Where did I leave the small gray bench?'', ``Which room should I go to?'',
``Is the gaming chair under the loft bed?''), while count/clock questions stayed flat
or improved.

Crucially, the flood is \emph{not} an IoU-threshold effect. A sweep of the 3D-bounds
merge threshold $t_{\text{iou}}\in\{0.0,0.15,0.3,0.45,0.6\}$ on six regressed qids left
the mean node count (213--233), the singleton fraction (0.48--0.52), and the answer
pattern (2--3/6) essentially unchanged. Repeated detections of the same object fail to
consolidate for structural reasons: YoloE class strings drift frame to frame (so the
label gate refuses a merge), mask-median centroids drift beyond the 0.42\,m spatial
gate, and partial-view 3D bounds overlap below any useful threshold. Tuning a single
IoU knob cannot fix this.

\subsection{Merge policy}
\label{app:merge_policy}

The graph-object-fusion config (\texttt{graph\_object\_fusion} block) therefore
separates the merge decision into ordered evidence gates plus admission, keep, and
growth controls:

\begin{description}
  \item[Gates] \texttt{gates.\{identity, bounds, embedding, spatial\}} order the
    evidence that can justify a merge, strongest first. Identity merges on a
    persistent track ID; bounds merge on 3D IoU; embedding merges on SigLIP-crop
    appearance similarity; spatial merges on centroid proximity (0.42\,m XY /
    0.55\,m 3D) with a nearest-node fallback radius. A candidate merges when the first
    active gate that matches passes.
  \item[Labels] \texttt{labels.\{require\_match, require\_match\_for\_instances\}}
    control when class strings gate a merge, with extra synonym / incompatible pairs
    for fixture dilution (e.g.\ \texttt{cab} $\leftrightarrow$ \texttt{cabinet}) and
    for pairs that must never merge (\texttt{person} vs \texttt{lamp}).
  \item[Keep] \texttt{keep.\{prefer\_support, union\_labels, union\_bounds,
    blend\_embedding, update\_xyz\}} decide which node survives a merge and how its
    anchor, bounds, labels, and embedding blend.
  \item[Admission] \texttt{admission.*} filters weak detections before they may
    enter the graph (confidence $>$ 0.12, $\geq$ 25 mask points).
  \item[Growth] \texttt{growth.\{max\_object\_nodes, temporal\_window\_steps\}} bound
    per-episode object-node count (0 = unlimited) and refuse merges into stale nodes.
\end{description}

\subsubsection{Appearance merging via SigLIP}

Instance detections initially carry no appearance embedding, so the embedding gate is
a no-op and label drift defeats merging. We attach a SigLIP crop embedding
(\texttt{gates.embedding.use\_siglip\_crops}) for every instance bbox; cosine
similarity $\geq$ 0.9 then permits a merge even when the label gate fails, while the
spatial gate still enforces proximity so two identical-but-distinct instances sitting
apart remain separate (preserving count recall).

\subsection{Lazy graph: close-look-only commits}
\label{app:lazy_graph}

An alternative to gate admission instead of merging: \texttt{LazyGraphController}
never streams YoloE detections into the scene graph during passive mapping. It records
lightweight viewpoint stamps for navigation, and runs the Qwen label extractor +
\texttt{add\_observation} only on a successful nav-arrival close look. Detector class
strings still feed voxel-level FIND recall but never author graph labels. This keeps
the graph at one node per deliberate inspection, which is the right regime for
location/state questions at the cost of denser per-instance count recall.

\subsection{Tying to precision}

The fp16-vs-int4 bake-off (Appendix~\ref{app:model_choice}) holds on both ingest
regimes: the precision lift is a model-weight effect, not an ingest artifact, and
subsets run under the same graph configuration remain internally comparable. We
recommend reporting the full-113 row under one frozen \texttt{graph\_object\_fusion}
block (see \texttt{configs/benchmarks/dynagraph.yaml} pin test) so a silent ingest
change can never re-baseline a published number again.
